% TOP_PROBLEM: {"id": 1365, "title": "Octal Games Periodicity", "category_id": 2, "category": {"id": 2, "name": "combinatorics", "display_name": "Combinatorics", "description": "Counting problems, graph theory, discrete structures.", "slug": "combinatorics", "order_index": 2, "created_at": "2026-07-31T15:26:25.671Z"}, "status": "open"}
% FIRST_POSTED: 2026-10-01
% ATTEMPT_STATUS: unresolved
% Imported from: unsolvedmath_additional_500.tex; UnsolvedMath ID 1365
% !TeX program = lualatex
% Compile twice: lualatex 1365.tex
% Self-contained: no external bibliography, images, or \input files.
% Numeric section labels and visible numbers are original UnsolvedMath IDs.
\documentclass[11pt,a4paper]{article}
\usepackage[margin=24mm,headheight=14pt]{geometry}
\usepackage{fontspec}
\setmainfont{Latin Modern Roman}
\setsansfont{Latin Modern Sans}
\setmonofont{Latin Modern Mono}
\usepackage{amsmath,amssymb,mathtools}
\usepackage{unicode-math}
\setmathfont{Latin Modern Math}
\usepackage{microtype}
\usepackage{enumitem,xurl,needspace}
\usepackage[dvipsnames]{xcolor}
\usepackage{hyperref}
\hypersetup{unicode=true,colorlinks=true,linkcolor=MidnightBlue,urlcolor=MidnightBlue,
 pdftitle={UnsolvedMath Problem 1365},pdfauthor={Source-annotated compilation},bookmarksnumbered=true}
\usepackage{bookmark,tocloft,titlesec,fancyhdr}
\setlength{\cftsecnumwidth}{4.2em}
\cftsetpnumwidth{2.5em}
\cftsetrmarg{4em}
\titleformat{\section}{\Large\bfseries\raggedright}{\thesection}{0.7em}{}
\pagestyle{fancy}\fancyhf{}
\fancyhead[L]{\small\sffamily Additional UnsolvedMath statements}
\fancyhead[R]{\small Original catalogue IDs}
\fancyfoot[C]{\thepage}
\setlength{\parindent}{0pt}
\setlength{\parskip}{5pt plus 1pt}
\setlength{\emergencystretch}{3em}
\setcounter{tocdepth}{1}\setcounter{secnumdepth}{1}
\setlist[itemize]{leftmargin=3em,itemsep=4pt,topsep=4pt,parsep=0pt}
\allowdisplaybreaks
\newcommand{\N}{\mathbb{N}}
\newcommand{\Z}{\mathbb{Z}}
\newcommand{\Q}{\mathbb{Q}}
\newcommand{\R}{\mathbb{R}}
\newcommand{\C}{\mathbb{C}}
\newcommand{\F}{\mathbb{F}}
\newcommand{\A}{\mathbb{A}}
\newcommand{\PP}{\mathbb{P}}
\newcommand{\E}{\mathbb{E}}
\newcommand{\eps}{\varepsilon}
\newcommand{\dd}{\,\mathrm d}
\newcommand{\sref}[2]{\hyperlink{source:#1:#2}{[S#2]}}
\newif\ifshowcatalogue
\showcataloguetrue
\newcommand{\cataloguescope}[1]{\ifshowcatalogue
 \paragraph{Statement recorded in the supplied source.}
 #1\fi}
\begin{document}
% UnsolvedMath ID 1365; source catalogue code GAME-008

\setcounter{section}{1364}

\section{Eventual Periodicity in Finite Octal Games}\label{1365}

{\small\sffamily UnsolvedMath ID 1365 · Combinatorics}

\subsection{Definitions and mathematical statement}

\paragraph{Shared notation.}Unless a section says otherwise, $\N=\{1,2,\ldots\}$,
$\Z$ is the ring of integers, $\Q,\R,\C$ are the rational, real, and complex
fields, $[N]=\{1,\ldots,\lfloor N\rfloor\}$, and $\log$ is the natural
logarithm. For real functions of a parameter tending to infinity,
$f\ll g$ and $f=O(g)$ mean $|f|\leq Cg$ eventually for some fixed $C$;
$f\gg g$ reverses this comparison; $f=o(g)$ means $f/g\to0$;
$f\sim g$ means $f/g\to1$; and $f\asymp g$ means both order bounds.
A subscript on $O$ or $\ll$ specifies allowed constant dependence.
The expression $n^{o(1)}$ in an upper bound means growth smaller than
$n^\epsilon$ for each fixed $\epsilon>0$. Local definitions govern any
exception, finite-field convention, or use of the same symbol for another
object. An asymptotic claim is not automatically an effective algorithm.



An octal take-and-break rule specifies, for each allowed number of removed counters, whether the remainder may consist of zero, one or two heaps; the three permissions are encoded by the bits of an octal digit. A finite code has only finitely many nonzero digits. Play is normal play on finite heaps. Its nim-sequence is $g(n)=\operatorname{mex}\{\bigoplus_i g(n_i):\text{legal move from a heap of size }n\}$, where $\oplus$ is binary xor. Eventual periodicity means $g(n+p)=g(n)$ for all $n\geq N_0$, for some $p,N_0$. \sref{1365}{1} \sref{1365}{2}

\paragraph{Statement recorded in the supplied source.}
Are the nim-sequences of all finite octal games eventually periodic? \sref{1365}{1}

\subsection{Short English statement}

Does the sequence of nim-values for every finitely specified octal heap game eventually repeat with a fixed period?

\subsection{Research attempt: a finite periodicity certificate and the nonsplitting subcase}
\textbf{Status.} We give a sufficient finite certificate for any proposed eventual period and prove eventual periodicity for every finite rule without splitting. We do not prove that a certificate exists for every finite octal code. The inspected primary data source [S2] uses bits for zero, one and two nonempty successor heaps and distinguishes empirical prefixes from proven periods.

\subsubsection{1. A safe finite certificate}
Let $K$ be the largest number of counters removed in a permitted move. A zero-removal move may split into two positive heaps, but may not leave the original heap unchanged. Take $N\geq1$, $p\geq1$, and put $T=2N+p+K$. We claim:
\[g(n+p)=g(n)\text{ for }N\leq n\leq T
\quad\Longrightarrow\quad g(n+p)=g(n)\text{ for every }n\geq N.\tag{1}\]
The bound is sufficient, not asserted optimal. Assume by induction that the equality is known below some $n>T$, and compare the option-value sets of heaps $n$ and $n+p$. A move removing the entire heap is unavailable because both sizes exceed $K$. A one-heap move removes $k\geq1$ counters, and its values agree by induction at base index $n-k\geq N$.

For a two-heap move from $n$, write $a+b=n-k$ with $a,b\geq1$. The sum exceeds $2N+p$, so one part, say $b$, is at least $N$. Replace it by $b+p$ to obtain a legal split of $n+p-k$ with the same xor value by induction. The comparison has base index $b<n$, even when $k=0$.

Conversely, in a split $a+b=n+p-k$ from $n+p$, the sum exceeds $2N+2p$, so one part is at least $N+p$. Reduce that part by $p$. It stays positive, the resulting split is legal from $n$, and the xor value is unchanged. Its comparison base index is at least $N$ and less than $n$. Octal rules permit equal parts whenever the two-heap bit is present, so equality between the reduced parts is not an extra prohibition. Both option sets are therefore equal, and applying mex proves the next equality in (1).

\subsubsection{2. An actual finite-state proof without splitting}
Suppose no two-heap bit occurs. With maximum removal $K\geq1$, there are at most $K$ one-heap options and at most one empty option. Thus every mex value is at most $K+1$. Once $n>K$, empty moves are impossible, and $g(n)$ is a fixed function of the previous $K$ values, determined by which removal digits permit one nonempty heap. The state vector of these $K$ values has at most $(K+2)^K$ possibilities. Iterating a deterministic map on this finite set eventually repeats a state and thereafter repeats forever. This proves eventual periodicity in the nonsplitting subcase. The zero rule $K=0$ is the constant-zero sequence.

This argument fails with splitting because a move can reference $g(a)$ for arbitrarily distant $a$, rather than just a bounded window behind $n$. Boundedness of values alone does not repair that missing finite-memory property.

\subsubsection{3. Certificate search and its logical limitation}
Given $N,p$, recurrence evaluation through $T+p=2N+2p+K$ checks the hypothesis of (1) by finite integer computations. Enumerating all pairs $N,p$ therefore semidecides eventual periodicity: if a period exists, one pair eventually passes and (1) certifies it. Failure of all pairs checked so far does not prove nonperiodicity. This is the same distinction as termination of a successful search versus termination on every input.

\subsubsection{4. Adversarial verification and remaining gap}
The diagnostic computes every two-digit code $0.d_1d_2$ through heap size 300, tests candidate certificates with $1\leq N\leq12$ and $1\leq p\leq12$, and checks all certificates found against the remaining computed tail. It also checks the nonsplitting codes against their finite-state cycles. These checks test indexing and the option correspondence, not the universal existence assertion. The proof explicitly treats zero-removal splitting, positive remainders and the disappearance of empty moves; an illegal one-heap zero-removal move would invalidate the well-founded recurrence and is excluded. What remains is a theorem forcing some $(N,p)$ to satisfy (1) for every finite splitting code, or a proof that a particular code never does.

\subsection{Sources}

{\small

\begin{itemize}

\item[\hypertarget{source:1365:1}{S1}] ULAM.AI, \emph{UnsolvedMath}, supplied \texttt{problems.json}, record 1365 (GAME-008), statement and background. The embedded literature-review date is 2026-08-17; that is the archive’s review date, not a fresh certification. Dataset landing page: \url{https://huggingface.co/datasets/ulamai/UnsolvedMath}.

\item[\hypertarget{source:1365:2}{S2}] A. Flammenkamp, Octal Games (). \url{https://wwwhomes.uni-bielefeld.de/achim/octal.html}. Bibliographic information imported from the supplied record.

\item[\hypertarget{source:1365:3}{S3}] Games of No Chance 5, open-problem list. \url{https://library.slmath.org/books/Book70/files/1005.pdf}. Bibliographic information imported from the supplied record.

\end{itemize}

}

\label{end:1365}
\end{document}
